% --- Archon physics formalization source begin ---
% archon:physics
% archon:covers PhyXMiniProblems/problem_phyx_mini_0911.lean
% archon:source-report reports/phyx_mini/problem_phyx_mini_0911.source.json

\chapter{Physics problem phyx\_mini\_0911}
\label{ch:PhyXMiniProblems_problem_phyx_mini_0911}

\paragraph{Problem source.}
\#\# Physical scenario

The water molecule \textbackslash{}(\textbackslash{}mathrm\{H\_2O\}\textbackslash{}) has a permanent dipole moment of magnitude \textbackslash{}(6.2\textbackslash{}times10\textasciicircum{}\{-30\}\textbackslash{},\textbackslash{}mathrm\{C\textbackslash{},m\}\textbackslash{}). A water molecule is located \textbackslash{}(10\textbackslash{},\textbackslash{}mathrm\{nm\}\textbackslash{}) from a \textbackslash{}(\textbackslash{}mathrm\{Na\textasciicircum{}+\}\textbackslash{}) ion in a saltwater solution.

\#\# Question

What force does the ion exert on the water molecule?

\#\# Answer choices

- A: A: 1.3 \textbackslash{}times 10\textasciicircum{}\{-14\}N

- B: B: 5.3 \textbackslash{}times 10\textasciicircum{}\{-14\}N

- C: C: 1.8 \textbackslash{}times 10\textasciicircum{}\{-14\}N

- D: D: 3.5 \textbackslash{}times 10\textasciicircum{}\{-14\}N

\#\# Figure caption (auxiliary; use the image as primary evidence)

The image illustrates the interaction between a sodium ion (\textbackslash{}( \textbackslash{}text\{Na\}\textasciicircum{}+ \textbackslash{})) and a water molecule. 

- On the left, there's a circle representing a sodium ion with a positive charge symbol inside. It's labeled "Na\textbackslash{}(\textasciicircum{}+\textbackslash{}) ion."
- On the right, there's a diagram representing a water molecule, with a partial negative charge (indicated by a minus sign) on one end and a partial positive charge (plus sign) on the other. It's labeled "Water molecule."
- Two arrows represent forces: 
  - An arrow labeled "\textbackslash{}(\textbackslash{}vec\{F\}\_\{\textbackslash{}text\{dipole on ion\}\}\textbackslash{})" points from the water molecule toward the sodium ion.
  - Another arrow labeled "\textbackslash{}(\textbackslash{}vec\{F\}\_\{\textbackslash{}text\{ion on dipole\}\}\textbackslash{})" points from the sodium ion toward the water molecule.
- A dashed arc between the two structures indicates the interaction, suggesting attraction.
- Below, there's a double-headed arrow showing the distance between the ion and the molecule, labeled \textbackslash{}( r = 10 \textbackslash{}, \textbackslash{}text\{nm\} \textbackslash{}).

The scene illustrates the concept of ion-dipole interactions.

\#\# Dataset metadata

Electrostatics; Physical Model Grounding Reasoning, Multi-Formula Reasoning

\paragraph{Recorded answer/context.}
C: C: 1.8 \textbackslash{}times 10\textasciicircum{}\{-14\}N

\paragraph{Figure/image path.}
/root/proposal\_for\_physic/science-mango/phyx\_mini\_run/phyx\_data/test\_image/911.png

\paragraph{Formalization target.}
create a compiling Lean file with sorry bodies at `PhyXMiniProblems/problem\_phyx\_mini\_0911.lean`. The Lean declarations must preserve the physical quantities, dimensions or dimensional roles, figure labels, governing-law hypotheses, and final relation expressed by this problem.
Use Mathlib/Physlib names found through LeanExplore where available. If a physics API is missing, introduce faithful local abstractions rather than scalar placeholder aliases.

\begin{theorem}[Physics formalization target]
\label{thm:physics:phyx_mini_0911:target}
\lean{PhyXMiniProblems.ProblemPhyXMini0911.problem_phyx_mini_0911}
\uses{def:physics:phyx-mini-0911:phyxminiproblems-problemphyxmini0911-forcecomponentinnewtons, def:physics:phyx-mini-0911:phyxminiproblems-problemphyxmini0911-sodiumionwaterdipolesetup, def:physics:phyx-mini-0911:phyxminiproblems-problemphyxmini0911-ionondipoleforcemagnitudeinnewtons, def:physics:phyx-mini-0911:phyxminiproblems-problemphyxmini0911-matchessodiumionwaterscenario, def:physics:phyx-mini-0911:phyxminiproblems-problemphyxmini0911-hasiondipolephysicaldata, def:physics:phyx-mini-0911:phyxminiproblems-problemphyxmini0911-matchessuppliediondipolefigure, def:physics:phyx-mini-0911:phyxminiproblems-problemphyxmini0911-hasionwateraxialgeometry, def:physics:phyx-mini-0911:phyxminiproblems-problemphyxmini0911-satisfiesalignediondipoleforcelaw, def:physics:phyx-mini-0911:phyxminiproblems-problemphyxmini0911-satisfiesiondipoleactionreaction, def:physics:phyx-mini-0911:phyxminiproblems-problemphyxmini0911-answerchoice, def:physics:phyx-mini-0911:phyxminiproblems-problemphyxmini0911-answerchoice-forceinnewtons, def:physics:phyx-mini-0911:phyxminiproblems-problemphyxmini0911-choicecdisplaytoleranceinnewtons}
The assigned autoformalize agent should translate this physics problem into Lean declarations in the covered file, with theorem and lemma proof bodies written as `by sorry`.
\end{theorem}
\begin{proof}
This is an autoformalization task, not a proof task. Produce statements that can later be proved by the physics prover without weakening the physical model.
\end{proof}

% --- Archon declaration topology begin ---
\section{Declaration topology}
\begin{definition}[force Dimension]
  \label{def:physics:phyx-mini-0911:phyxminiproblems-problemphyxmini0911-forcedimension}
  \lean{PhyXMiniProblems.ProblemPhyXMini0911.forceDimension}
  The physical dimension M L T⁻² of force
\end{definition}
\begin{proof}
  This is the declared model definition of force Dimension.
\end{proof}
\begin{definition}[electric Dipole Moment Dimension]
  \label{def:physics:phyx-mini-0911:phyxminiproblems-problemphyxmini0911-electricdipolemomentdimension}
  \lean{PhyXMiniProblems.ProblemPhyXMini0911.electricDipoleMomentDimension}
  The physical dimension C L of electric dipole moment
\end{definition}
\begin{proof}
  This is the declared model definition of electric Dipole Moment Dimension.
\end{proof}
\begin{definition}[Signed Charge Quantity]
  \label{def:physics:phyx-mini-0911:phyxminiproblems-problemphyxmini0911-signedchargequantity}
  \lean{PhyXMiniProblems.ProblemPhyXMini0911.SignedChargeQuantity}
  A signed, unit-independent electric charge
\end{definition}
\begin{proof}
  This is the declared model definition of Signed Charge Quantity.
\end{proof}
\begin{definition}[Dipole Moment Magnitude Quantity]
  \label{def:physics:phyx-mini-0911:phyxminiproblems-problemphyxmini0911-dipolemomentmagnitudequantity}
  \lean{PhyXMiniProblems.ProblemPhyXMini0911.DipoleMomentMagnitudeQuantity}
  \uses{def:physics:phyx-mini-0911:phyxminiproblems-problemphyxmini0911-electricdipolemomentdimension}
  A nonnegative, unit-independent electric dipole-moment magnitude
\end{definition}
\begin{proof}
  This is the declared model definition of Dipole Moment Magnitude Quantity.
\end{proof}
\begin{definition}[Axial Position Quantity]
  \label{def:physics:phyx-mini-0911:phyxminiproblems-problemphyxmini0911-axialpositionquantity}
  \lean{PhyXMiniProblems.ProblemPhyXMini0911.AxialPositionQuantity}
  A signed axial position
\end{definition}
\begin{proof}
  This is the declared model definition of Axial Position Quantity.
\end{proof}
\begin{definition}[Length Quantity]
  \label{def:physics:phyx-mini-0911:phyxminiproblems-problemphyxmini0911-lengthquantity}
  \lean{PhyXMiniProblems.ProblemPhyXMini0911.LengthQuantity}
  A nonnegative, unit-independent physical separation
\end{definition}
\begin{proof}
  This is the declared model definition of Length Quantity.
\end{proof}
\begin{definition}[Axial Force Component Quantity]
  \label{def:physics:phyx-mini-0911:phyxminiproblems-problemphyxmini0911-axialforcecomponentquantity}
  \lean{PhyXMiniProblems.ProblemPhyXMini0911.AxialForceComponentQuantity}
  \uses{def:physics:phyx-mini-0911:phyxminiproblems-problemphyxmini0911-forcedimension}
  A signed axial component of a physical force
\end{definition}
\begin{proof}
  This is the declared model definition of Axial Force Component Quantity.
\end{proof}
\begin{definition}[charge In Coulombs]
  \label{def:physics:phyx-mini-0911:phyxminiproblems-problemphyxmini0911-chargeincoulombs}
  \lean{PhyXMiniProblems.ProblemPhyXMini0911.chargeInCoulombs}
  \uses{def:physics:phyx-mini-0911:phyxminiproblems-problemphyxmini0911-signedchargequantity}
  Read a signed physical charge in coherent-SI coulombs
\end{definition}
\begin{proof}
  This is the declared model definition of charge In Coulombs.
\end{proof}
\begin{definition}[dipole Moment In Coulomb Meters]
  \label{def:physics:phyx-mini-0911:phyxminiproblems-problemphyxmini0911-dipolemomentincoulombmeters}
  \lean{PhyXMiniProblems.ProblemPhyXMini0911.dipoleMomentInCoulombMeters}
  \uses{def:physics:phyx-mini-0911:phyxminiproblems-problemphyxmini0911-dipolemomentmagnitudequantity}
  Read a dipole-moment magnitude in coherent-SI coulomb-metres
\end{definition}
\begin{proof}
  This is the declared model definition of dipole Moment In Coulomb Meters.
\end{proof}
\begin{definition}[position In Meters]
  \label{def:physics:phyx-mini-0911:phyxminiproblems-problemphyxmini0911-positioninmeters}
  \lean{PhyXMiniProblems.ProblemPhyXMini0911.positionInMeters}
  \uses{def:physics:phyx-mini-0911:phyxminiproblems-problemphyxmini0911-axialpositionquantity}
  Read an axial position in coherent-SI metres
\end{definition}
\begin{proof}
  This is the declared model definition of position In Meters.
\end{proof}
\begin{definition}[length In Meters]
  \label{def:physics:phyx-mini-0911:phyxminiproblems-problemphyxmini0911-lengthinmeters}
  \lean{PhyXMiniProblems.ProblemPhyXMini0911.lengthInMeters}
  \uses{def:physics:phyx-mini-0911:phyxminiproblems-problemphyxmini0911-lengthquantity}
  Read a nonnegative physical length in coherent-SI metres
\end{definition}
\begin{proof}
  This is the declared model definition of length In Meters.
\end{proof}
\begin{definition}[length In Nanometers]
  \label{def:physics:phyx-mini-0911:phyxminiproblems-problemphyxmini0911-lengthinnanometers}
  \lean{PhyXMiniProblems.ProblemPhyXMini0911.lengthInNanometers}
  \uses{def:physics:phyx-mini-0911:phyxminiproblems-problemphyxmini0911-lengthquantity, def:physics:phyx-mini-0911:phyxminiproblems-problemphyxmini0911-lengthinmeters}
  Read a physical length in nanometres, the unit printed in the image
\end{definition}
\begin{proof}
  This is the declared model definition of length In Nanometers.
\end{proof}
\begin{definition}[force Component In Newtons]
  \label{def:physics:phyx-mini-0911:phyxminiproblems-problemphyxmini0911-forcecomponentinnewtons}
  \lean{PhyXMiniProblems.ProblemPhyXMini0911.forceComponentInNewtons}
  \uses{def:physics:phyx-mini-0911:phyxminiproblems-problemphyxmini0911-axialforcecomponentquantity}
  Read a signed axial force component in coherent-SI newtons
\end{definition}
\begin{proof}
  This is the declared model definition of force Component In Newtons.
\end{proof}
\begin{definition}[Figure Object]
  \label{def:physics:phyx-mini-0911:phyxminiproblems-problemphyxmini0911-figureobject}
  \lean{PhyXMiniProblems.ProblemPhyXMini0911.FigureObject}
  The two physical objects named in image 911.png
\end{definition}
\begin{proof}
  This is the declared model definition of Figure Object.
\end{proof}
\begin{definition}[Water Molecule End]
  \label{def:physics:phyx-mini-0911:phyxminiproblems-problemphyxmini0911-watermoleculeend}
  \lean{PhyXMiniProblems.ProblemPhyXMini0911.WaterMoleculeEnd}
  The two ends of the pictured water-dipole representation
\end{definition}
\begin{proof}
  This is the declared model definition of Water Molecule End.
\end{proof}
\begin{definition}[Partial Charge Sign]
  \label{def:physics:phyx-mini-0911:phyxminiproblems-problemphyxmini0911-partialchargesign}
  \lean{PhyXMiniProblems.ProblemPhyXMini0911.PartialChargeSign}
  The partial-charge sign drawn at an end of the water molecule
\end{definition}
\begin{proof}
  This is the declared model definition of Partial Charge Sign.
\end{proof}
\begin{definition}[Force Arrow]
  \label{def:physics:phyx-mini-0911:phyxminiproblems-problemphyxmini0911-forcearrow}
  \lean{PhyXMiniProblems.ProblemPhyXMini0911.ForceArrow}
  The two force arrows explicitly labelled in the primary image
\end{definition}
\begin{proof}
  This is the declared model definition of Force Arrow.
\end{proof}
\begin{definition}[Axial Direction]
  \label{def:physics:phyx-mini-0911:phyxminiproblems-problemphyxmini0911-axialdirection}
  \lean{PhyXMiniProblems.ProblemPhyXMini0911.AxialDirection}
  Direction on the image's horizontal axis, with positive coordinate rightward
\end{definition}
\begin{proof}
  This is the declared model definition of Axial Direction.
\end{proof}
\begin{definition}[Interaction Medium]
  \label{def:physics:phyx-mini-0911:phyxminiproblems-problemphyxmini0911-interactionmedium}
  \lean{PhyXMiniProblems.ProblemPhyXMini0911.InteractionMedium}
  The material environment stated in the prose
\end{definition}
\begin{proof}
  This is the declared model definition of Interaction Medium.
\end{proof}
\begin{definition}[Ion Model]
  \label{def:physics:phyx-mini-0911:phyxminiproblems-problemphyxmini0911-ionmodel}
  \lean{PhyXMiniProblems.ProblemPhyXMini0911.IonModel}
  Idealization used for the sodium ion in the textbook computation
\end{definition}
\begin{proof}
  This is the declared model definition of Ion Model.
\end{proof}
\begin{definition}[Dipole Orientation]
  \label{def:physics:phyx-mini-0911:phyxminiproblems-problemphyxmini0911-dipoleorientation}
  \lean{PhyXMiniProblems.ProblemPhyXMini0911.DipoleOrientation}
  Orientation of the permanent water dipole relative to the positive ion
\end{definition}
\begin{proof}
  This is the declared model definition of Dipole Orientation.
\end{proof}
\begin{definition}[Electrostatic Approximation]
  \label{def:physics:phyx-mini-0911:phyxminiproblems-problemphyxmini0911-electrostaticapproximation}
  \lean{PhyXMiniProblems.ProblemPhyXMini0911.ElectrostaticApproximation}
  The recorded numerical answer uses the unscreened textbook ion--dipole law; this constructor makes that approximation explicit despite the stated saltwater environment
\end{definition}
\begin{proof}
  This is the declared model definition of Electrostatic Approximation.
\end{proof}
\begin{definition}[expected Object Label]
  \label{def:physics:phyx-mini-0911:phyxminiproblems-problemphyxmini0911-expectedobjectlabel}
  \lean{PhyXMiniProblems.ProblemPhyXMini0911.expectedObjectLabel}
  \uses{def:physics:phyx-mini-0911:phyxminiproblems-problemphyxmini0911-figureobject}
  Literal object labels visible below the two drawings
\end{definition}
\begin{proof}
  This is the declared model definition of expected Object Label.
\end{proof}
\begin{definition}[expected Water End Sign]
  \label{def:physics:phyx-mini-0911:phyxminiproblems-problemphyxmini0911-expectedwaterendsign}
  \lean{PhyXMiniProblems.ProblemPhyXMini0911.expectedWaterEndSign}
  \uses{def:physics:phyx-mini-0911:phyxminiproblems-problemphyxmini0911-watermoleculeend, def:physics:phyx-mini-0911:phyxminiproblems-problemphyxmini0911-partialchargesign}
  Partial-charge signs read from the two ends of the water molecule
\end{definition}
\begin{proof}
  This is the declared model definition of expected Water End Sign.
\end{proof}
\begin{definition}[expected Force Arrow Label]
  \label{def:physics:phyx-mini-0911:phyxminiproblems-problemphyxmini0911-expectedforcearrowlabel}
  \lean{PhyXMiniProblems.ProblemPhyXMini0911.expectedForceArrowLabel}
  \uses{def:physics:phyx-mini-0911:phyxminiproblems-problemphyxmini0911-forcearrow}
  Literal names of the force arrows in the figure
\end{definition}
\begin{proof}
  This is the declared model definition of expected Force Arrow Label.
\end{proof}
\begin{definition}[expected Force Arrow Direction]
  \label{def:physics:phyx-mini-0911:phyxminiproblems-problemphyxmini0911-expectedforcearrowdirection}
  \lean{PhyXMiniProblems.ProblemPhyXMini0911.expectedForceArrowDirection}
  \uses{def:physics:phyx-mini-0911:phyxminiproblems-problemphyxmini0911-forcearrow, def:physics:phyx-mini-0911:phyxminiproblems-problemphyxmini0911-axialdirection}
  The image shows both interaction forces pointing toward the other object
\end{definition}
\begin{proof}
  This is the declared model definition of expected Force Arrow Direction.
\end{proof}
\begin{definition}[Sodium Ion Water Figure]
  \label{def:physics:phyx-mini-0911:phyxminiproblems-problemphyxmini0911-sodiumionwaterfigure}
  \lean{PhyXMiniProblems.ProblemPhyXMini0911.SodiumIonWaterFigure}
  \uses{def:physics:phyx-mini-0911:phyxminiproblems-problemphyxmini0911-figureobject, def:physics:phyx-mini-0911:phyxminiproblems-problemphyxmini0911-watermoleculeend, def:physics:phyx-mini-0911:phyxminiproblems-problemphyxmini0911-partialchargesign, def:physics:phyx-mini-0911:phyxminiproblems-problemphyxmini0911-forcearrow, def:physics:phyx-mini-0911:phyxminiproblems-problemphyxmini0911-axialdirection}
  Presentation data transcribed from the primary bitmap. In particular, this structure contains no numerical physical-force value
\end{definition}
\begin{proof}
  This is the declared model definition of Sodium Ion Water Figure.
\end{proof}
\begin{definition}[Sodium Ion Water Dipole Setup]
  \label{def:physics:phyx-mini-0911:phyxminiproblems-problemphyxmini0911-sodiumionwaterdipolesetup}
  \lean{PhyXMiniProblems.ProblemPhyXMini0911.SodiumIonWaterDipoleSetup}
  \uses{def:physics:phyx-mini-0911:phyxminiproblems-problemphyxmini0911-signedchargequantity, def:physics:phyx-mini-0911:phyxminiproblems-problemphyxmini0911-dipolemomentmagnitudequantity, def:physics:phyx-mini-0911:phyxminiproblems-problemphyxmini0911-axialpositionquantity, def:physics:phyx-mini-0911:phyxminiproblems-problemphyxmini0911-lengthquantity, def:physics:phyx-mini-0911:phyxminiproblems-problemphyxmini0911-axialforcecomponentquantity, def:physics:phyx-mini-0911:phyxminiproblems-problemphyxmini0911-interactionmedium, def:physics:phyx-mini-0911:phyxminiproblems-problemphyxmini0911-ionmodel, def:physics:phyx-mini-0911:phyxminiproblems-problemphyxmini0911-dipoleorientation, def:physics:phyx-mini-0911:phyxminiproblems-problemphyxmini0911-electrostaticapproximation, def:physics:phyx-mini-0911:phyxminiproblems-problemphyxmini0911-sodiumionwaterfigure}
  Independent physical quantities associated with the prose and diagram. The two force fields are observables, and the remainder is an independent model discrepancy; none is a definition of the requested answer
\end{definition}
\begin{proof}
  This is the declared model definition of Sodium Ion Water Dipole Setup.
\end{proof}
\begin{definition}[ion On Dipole Force Magnitude In Newtons]
  \label{def:physics:phyx-mini-0911:phyxminiproblems-problemphyxmini0911-ionondipoleforcemagnitudeinnewtons}
  \lean{PhyXMiniProblems.ProblemPhyXMini0911.ionOnDipoleForceMagnitudeInNewtons}
  \uses{def:physics:phyx-mini-0911:phyxminiproblems-problemphyxmini0911-forcecomponentinnewtons, def:physics:phyx-mini-0911:phyxminiproblems-problemphyxmini0911-sodiumionwaterdipolesetup}
  Magnitude of the requested force, obtained from its signed axial component
\end{definition}
\begin{proof}
  This is the declared model definition of ion On Dipole Force Magnitude In Newtons.
\end{proof}
\begin{definition}[Matches Sodium Ion Water Scenario]
  \label{def:physics:phyx-mini-0911:phyxminiproblems-problemphyxmini0911-matchessodiumionwaterscenario}
  \lean{PhyXMiniProblems.ProblemPhyXMini0911.MatchesSodiumIonWaterScenario}
  \uses{def:physics:phyx-mini-0911:phyxminiproblems-problemphyxmini0911-chargeincoulombs, def:physics:phyx-mini-0911:phyxminiproblems-problemphyxmini0911-sodiumionwaterdipolesetup}
  Qualitative physical modeling assumptions stated or implied by the scenario
\end{definition}
\begin{proof}
  This is the declared model definition of Matches Sodium Ion Water Scenario.
\end{proof}
\begin{definition}[Has Ion Dipole Physical Data]
  \label{def:physics:phyx-mini-0911:phyxminiproblems-problemphyxmini0911-hasiondipolephysicaldata}
  \lean{PhyXMiniProblems.ProblemPhyXMini0911.HasIonDipolePhysicalData}
  \uses{def:physics:phyx-mini-0911:phyxminiproblems-problemphyxmini0911-chargeincoulombs, def:physics:phyx-mini-0911:phyxminiproblems-problemphyxmini0911-dipolemomentincoulombmeters, def:physics:phyx-mini-0911:phyxminiproblems-problemphyxmini0911-sodiumionwaterdipolesetup}
  Numerical physical data from the prose together with the standard elementary charge and Coulomb-constant values needed by the textbook model. No force value from an answer choice appears here
\end{definition}
\begin{proof}
  This is the declared model definition of Has Ion Dipole Physical Data.
\end{proof}
\begin{definition}[Matches Supplied Ion Dipole Figure]
  \label{def:physics:phyx-mini-0911:phyxminiproblems-problemphyxmini0911-matchessuppliediondipolefigure}
  \lean{PhyXMiniProblems.ProblemPhyXMini0911.MatchesSuppliedIonDipoleFigure}
  \uses{def:physics:phyx-mini-0911:phyxminiproblems-problemphyxmini0911-lengthinnanometers, def:physics:phyx-mini-0911:phyxminiproblems-problemphyxmini0911-expectedobjectlabel, def:physics:phyx-mini-0911:phyxminiproblems-problemphyxmini0911-expectedwaterendsign, def:physics:phyx-mini-0911:phyxminiproblems-problemphyxmini0911-expectedforcearrowlabel, def:physics:phyx-mini-0911:phyxminiproblems-problemphyxmini0911-expectedforcearrowdirection, def:physics:phyx-mini-0911:phyxminiproblems-problemphyxmini0911-sodiumionwaterdipolesetup}
  All qualitative and scalar evidence read from image 911.png, including the calibration between the printed distance and the physical separation
\end{definition}
\begin{proof}
  This is the declared model definition of Matches Supplied Ion Dipole Figure.
\end{proof}
\begin{definition}[Has Ion Water Axial Geometry]
  \label{def:physics:phyx-mini-0911:phyxminiproblems-problemphyxmini0911-hasionwateraxialgeometry}
  \lean{PhyXMiniProblems.ProblemPhyXMini0911.HasIonWaterAxialGeometry}
  \uses{def:physics:phyx-mini-0911:phyxminiproblems-problemphyxmini0911-positioninmeters, def:physics:phyx-mini-0911:phyxminiproblems-problemphyxmini0911-lengthinmeters, def:physics:phyx-mini-0911:phyxminiproblems-problemphyxmini0911-sodiumionwaterdipolesetup}
  The signed coordinates realize the left-to-right ion--water order and the positive separation shown by the double-headed arrow
\end{definition}
\begin{proof}
  This is the declared model definition of Has Ion Water Axial Geometry.
\end{proof}
\begin{definition}[Satisfies Aligned Ion Dipole Force Law]
  \label{def:physics:phyx-mini-0911:phyxminiproblems-problemphyxmini0911-satisfiesalignediondipoleforcelaw}
  \lean{PhyXMiniProblems.ProblemPhyXMini0911.SatisfiesAlignedIonDipoleForceLaw}
  \uses{def:physics:phyx-mini-0911:phyxminiproblems-problemphyxmini0911-chargeincoulombs, def:physics:phyx-mini-0911:phyxminiproblems-problemphyxmini0911-dipolemomentincoulombmeters, def:physics:phyx-mini-0911:phyxminiproblems-problemphyxmini0911-lengthinmeters, def:physics:phyx-mini-0911:phyxminiproblems-problemphyxmini0911-forcecomponentinnewtons, def:physics:phyx-mini-0911:phyxminiproblems-problemphyxmini0911-sodiumionwaterdipolesetup}
  For a point ion and an aligned permanent dipole on the same axis, the leading force magnitude is 2 k |q| p / r\textasciicircum{}3. With positive coordinate rightward and the positive ion on the left, the negative sign states that the leading force on the water dipole is attractive. The water molecule is not literally a mathematical point dipole, and the scenario says that the interaction occurs in saltwater. Consequently the textbook formula is not asserted as an exact physical identity.
\end{definition}
\begin{proof}
  This is the declared model definition of Satisfies Aligned Ion Dipole Force Law.
\end{proof}
\begin{definition}[Satisfies Ion Dipole Action Reaction]
  \label{def:physics:phyx-mini-0911:phyxminiproblems-problemphyxmini0911-satisfiesiondipoleactionreaction}
  \lean{PhyXMiniProblems.ProblemPhyXMini0911.SatisfiesIonDipoleActionReaction}
  \uses{def:physics:phyx-mini-0911:phyxminiproblems-problemphyxmini0911-forcecomponentinnewtons, def:physics:phyx-mini-0911:phyxminiproblems-problemphyxmini0911-sodiumionwaterdipolesetup}
  Newton's third law relates the two force arrows shown in the image
\end{definition}
\begin{proof}
  This is the declared model definition of Satisfies Ion Dipole Action Reaction.
\end{proof}
\begin{definition}[Answer Choice]
  \label{def:physics:phyx-mini-0911:phyxminiproblems-problemphyxmini0911-answerchoice}
  \lean{PhyXMiniProblems.ProblemPhyXMini0911.AnswerChoice}
  Labels of the four force magnitudes printed in the problem source
\end{definition}
\begin{proof}
  This is the declared model definition of Answer Choice.
\end{proof}
\begin{definition}[force In Newtons]
  \label{def:physics:phyx-mini-0911:phyxminiproblems-problemphyxmini0911-answerchoice-forceinnewtons}
  \lean{PhyXMiniProblems.ProblemPhyXMini0911.AnswerChoice.forceInNewtons}
  \uses{def:physics:phyx-mini-0911:phyxminiproblems-problemphyxmini0911-answerchoice}
  Force magnitude in newtons printed beside each answer label
\end{definition}
\begin{proof}
  This is the declared model definition of force In Newtons.
\end{proof}
\begin{definition}[choice CDisplay Tolerance In Newtons]
  \label{def:physics:phyx-mini-0911:phyxminiproblems-problemphyxmini0911-choicecdisplaytoleranceinnewtons}
  \lean{PhyXMiniProblems.ProblemPhyXMini0911.choiceCDisplayToleranceInNewtons}
  Half of the last displayed decimal place in answer choice C
\end{definition}
\begin{proof}
  This is the declared model definition of choice CDisplay Tolerance In Newtons.
\end{proof}
\begin{definition}[recorded Dataset Answer]
  \label{def:physics:phyx-mini-0911:phyxminiproblems-problemphyxmini0911-recordeddatasetanswer}
  \lean{PhyXMiniProblems.ProblemPhyXMini0911.recordedDatasetAnswer}
  \uses{def:physics:phyx-mini-0911:phyxminiproblems-problemphyxmini0911-answerchoice}
  The answer label recorded by the supplied dataset
\end{definition}
\begin{proof}
  This is the declared model definition of recorded Dataset Answer.
\end{proof}
% --- Archon declaration topology end ---
% --- Archon physics formalization source end ---
